\title{Controlling Text-to-Image Diffusion by Orthogonal Finetuning}
\subsection{Efficient Orthogonal Parameterization}\label{sect:eff_ortho}

\textbf{Balancing flexibility and regularization in finetuning}. Our investigation highlights a fundamental balance between flexibility and regularity. While conventional finetuning affords maximum flexibility, it tends to lack stability and is prone to model collapse. OFT, despite its simplicity, effectively negotiates this balance by maintaining the pairwise angles between neurons. Moreover, the block-diagonal approach imposes a stricter regularization constraint on the orthogonal transformation.

In response to the first question, we build upon previous empirical findings presented in \cite{liu2018decoupled,chen2020angular}, which indicate that angular differences in features effectively capture semantic disparities. SphereNet~\cite{liu2017deep} demonstrates that restricting all neuron activations to lie on a unit hypersphere during training allows a neural network to retain its representational power and often results in enhanced generalization performance. This suggests that the informational content of neural responses is predominantly encoded in their direction rather than magnitude. To illustrate the relevance of neuron angles, we designed a controlled experiment (Figure~\ref{fig:toy_ae}) by training a standard convolutional autoencoder on flower images. During training, feature maps were generated via the conventional inner product ($z$ is the elementwise output from convolution kernel $\bm{w}$ and input $\bm{x}$ within a sliding window). In evaluation, we contrasted three approaches for producing feature maps: (a) replicating the inner product from training, (b) isolating only magnitude, and (c) extracting solely angular (directional) information. As displayed in Figure~\ref{fig:toy_ae}, reconstructing images using just the neurons' angular component achieves nearly perfect fidelity to the original inputs, whereas magnitude alone carries no meaningful information. Notably, the cosine-based activation (c) was never applied during the training phase—the training depended exclusively on the inner product. These outcomes underscore the pivotal role of neuron direction in preserving the semantics of input data, indicating that adjusting neuron orientation is likely the most effective strategy for altering image semantics.

Here, the orthogonality of $\bm{R}$ is preserved while an additional $\epsilon$-ball constraint enforces proximity to the identity. The Cayley transform (see Section~\ref{sect:eff_ortho}) allows parameterization of orthogonality; however, enforcing the $\epsilon$-deviation constraint is challenging with this formulation. To shed light on this, we approximate the Cayley-transformed matrix $\thickmuskip=2mu \medmuskip=2mu \bm{R}=(\bm{I}+\bm{Q})(I-\bm{Q})^{-1}$ using a truncated Neumann series, yielding $\thickmuskip=2mu \medmuskip=2mu \bm{R}\approx\bm{I}+2\bm{Q}+\mathcal{O}(\bm{Q}^2)$, assuming the expansion is valid.

where $\text{O}(d)$ represents the group of $d$-dimensional orthogonal matrices, with $\thickmuskip=2mu \medmuskip=2mu \bm{R}\in\mathbb{R}^{d\times d}$ and $\thickmuskip=2mu \medmuskip=2mu \bm{R}_i\in\mathbb{R}^{d/r\times d/r}$ for all $i$ denoting orthogonal matrices. If $\thickmuskip=2mu \medmuskip=2mu r=1$, the block-diagonal structure collapses to a fully general orthogonal matrix. The number of independent parameters in a $d\times d$ orthogonal matrix is $\thickmuskip=2mu \medmuskip=2mu d(d-1)/2$, leading to a parameter complexity of $\mathcal{O}(d^2)$. For an orthogonal matrix composed of $r$ blocks of size $d/r \times d/r$, the total number of parameters reduces to $\thickmuskip=2mu \medmuskip=2mu d(d/r-1)/2$, corresponding to complexity $\thickmuskip=2mu \medmuskip=2mu \mathcal{O}(d^2/r)$. An optional further reduction is possible by enforcing block sharing, that is, by setting $\thickmuskip=2mu \medmuskip=2mu\bm{R}_i=\bm{R}_j$ for all $i\neq j$, which brings the number of parameters down to $\mathcal{O}(d^2/r^2)$. Notwithstanding these parameter reduction techniques, the global matrix $\bm{R}$ maintains its orthogonality, ensuring hyperspherical energy is fully retained.

\begin{itemize}[leftmargin=*,nosep]

Here, $\thickmuskip=2mu \medmuskip=2mu \hat{\bm{w}}_i:=\bm{w}_i/\|\bm{w}_i\|$ represents the normalized version of the $i$th neuron, and the hyperspherical energy $\text{HE}(\bm{W})$ for a fully connected layer is defined as $\thickmuskip=2mu \medmuskip=2mu \sum_{i\neq j} \|\hat{\bm{w}}_i-\hat{\bm{w}}_j\|^{-1}$. It is clear that the expression in Eq.~\eqref{eq:oft_principle} is minimized (to zero) when $\bm{W}$ and $\bm{W}^0$ are related by an orthogonal transformation, i.e., $\thickmuskip=2mu \medmuskip=2mu \bm{W}=\bm{R}\bm{W}^0$, where $\thickmuskip=2mu \medmuskip=2mu \bm{R}\in\mathbb{R}^{d\times d}$ is an orthogonal matrix—its determinant indicating whether the transformation is a rotation ($1$) or a reflection ($-1$). This insight underpins OFT: the key is to adjust the neural network only by learning layer-wise shared orthogonal matrices that transform the neurons in each layer. More

\textbf{Quantitative comparison}. To assess the efficacy of controllable generation, we propose a control consistency metric. This metric involves extracting the control signal from each synthesized image and contrasting it with the corresponding original input control signal. Specifically, for the C2I task, we report IoU and F1 scores. For the S2I task, metrics include mean IoU, mean accuracy, and overall accuracy. Regarding the L2F task, the evaluation consists of computing the average $\ell_2$ distance between predicted and input control landmarks. Further information about these consistency metrics is provided in Appendix~\ref{app:settings}. Each method is assessed with its optimal set of hyperparameters. The findings summarized in Table~\ref{table:control_gen} indicate that OFT and COFT demonstrate notably superior and more precise control compared to alternative approaches. Adapter-based techniques (\emph{e.g.}, T2I-Adapter and ControlNet) display slower convergence and ultimately underperform relative to the leading methods. LoRA outperforms ControlNet for the S2I task but lags behind on C2I and L2F tasks. Overall, the maximal performance of LoRA remains modest, even with thorough hyperparameter optimization. By contrast, OFT shows no clear saturation in performance; increasing the count of trainable parameters continues to yield gains empirically. It is worth highlighting that our quantitative evaluation strategy for controllable generation constitutes a novel aspect of our work, as it enables precise assessment of the control capabilities in finetuned models (to the extent allowed by the accuracy of the employed segmentation/detection models).

\section{Experiments and Results}


\begin{abstract}
Large-scale text-to-image diffusion models have demonstrated remarkable proficiency in synthesizing photorealistic images from textual descriptions. However, devising effective strategies to steer or constrain these advanced models for diverse downstream applications remains a significant and unresolved issue. In response, we present a theoretically grounded adaptation technique—Orthogonal Finetuning (OFT)—for tailoring text-to-image diffusion models to new tasks. Distinct from existing approaches, OFT is theoretically guaranteed to maintain the hyperspherical energy, which encapsulates the pairwise interactions among neurons residing on the unit hypersphere. Our analysis suggests that preserving this geometric property is vital for retaining the semantic generative capability intrinsic to text-to-image diffusion models. Furthermore, to enhance the robustness of finetuning, we introduce Constrained Orthogonal Finetuning (COFT), which incorporates a radius constraint on the hypersphere. We focus on two key finetuning paradigms for text-to-image models: subject-driven generation, where the objective is to synthesize images of a specific subject based on a limited set of reference images and a textual prompt; and controllable generation, where the model is further conditioned on auxiliary signals. Through empirical evaluation, we demonstrate that our OFT framework consistently outperforms prior methods in both generative fidelity and optimization efficiency.
\end{abstract}

\subsection{General Framework}

\begin{equation}
    \footnotesize
    \bm{R}=\text{diag}(\bm{R}_1,\bm{R}_2,\cdots,\bm{R}_r)=\begin{bmatrix}
    \bm{R}_1\in \text{O}(\frac{d}{r}) &  &\\
    &\ddots & \\
    & & \bm{R}_r\in \text{O}(\frac{d}{r})
    \end{bmatrix}\in \text{O}(d)
\end{equation}

Since the series converges in the operator norm, it is permissible to transfer the constraint $\thickmuskip=2mu \medmuskip=2mu\|\bm{R}-\bm{I}\|\leq\epsilon$ directly into the domain of the Cayley transform. This reformulates the original constraint as $\thickmuskip=2mu \medmuskip=2mu \|\bm{Q}-\bm{0}\|\leq\epsilon'$, where $\bm{0}$ represents a matrix of all zeros and $\epsilon'$ is a distinct tolerance parameter from $\epsilon$. The imposed restriction on $\bm{Q}$ can be straightforwardly maintained using projected gradient descent. To initialize the orthogonal matrix $\bm{R}$ as the identity, we simply set $\bm{Q}$ to the zero matrix. From this perspective, COFT imposes two concrete requirements: limiting changes in hyperspherical energy and strictly bounding how far the finetuned parameters drift from those of the pretrained network. Although the second requirement often remains implicit in prior finetuning strategies, COFT explicitly encodes it. Our empirical observations suggest that this explicit deviation constraint further enhances the stability of COFT compared to OFT. Figure~\ref{fig:eps_coft} illustrates the influence of varying $\epsilon$ on COFT performance. Larger $\epsilon$ grants more freedom in finetuning and yields results close to standard OFT, whereas smaller $\epsilon$ restricts COFT's changes, rendering its output closer to that of the pretrained text-to-image diffusion baseline.

Nevertheless, OFT gives rise to several intriguing unresolved questions. Firstly, the method enforces orthogonality using the Cayley parametrization, which necessitates computing a matrix inverse, thereby somewhat constraining OFT's scalability. While a block diagonal parametrization is introduced to partially mitigate this limitation, discovering differentiable and efficient methods for performing matrix inversion remains an unresolved technical obstacle. Secondly, OFT exhibits promising compositional properties: orthogonal matrices resulting from distinct OFT finetuning tasks can be sequentially multiplied, yielding another orthogonal matrix. Investigating whether these composed matrices can collectively retain knowledge from all underlying downstream tasks presents a fertile area for future research. Lastly, OFT’s parameter efficiency is intimately tied to the block diagonal construction, which introduces certain biases and can limit flexibility. Developing more effective and less biased strategies to enhance parameter efficiency is a crucial open challenge.

Here, $\bm{W}$ indicates the OFT-finetuned weight matrix, and $\bm{I}$ denotes the identity matrix. Figure~\ref{fig:block} provides an illustration of OFT. Analogous to the zero initialization strategy employed in LoRA, OFT must be configured to ensure that fine-tuning proceeds exactly from $\bm{W}^0$. This is accomplished by setting the initial value of the orthogonal matrix $\bm{R}$ to the identity matrix, guaranteeing that the fine-tuned model begins with the pretrained parameters. To maintain the orthogonality of $\bm{R}$ during optimization, one can employ differentiable orthogonalization techniques as described in \cite{liu2021orthogonal,lezcano2019cheap}. In what follows, we examine how orthogonality can be preserved in an efficient manner.

The rationale behind employing orthogonal transformations for neuron finetuning is partly motivated by insights from the 2D Fourier transform, where an image is represented by its magnitude and phase spectra. It is well established that the phase spectrum—capturing the angular relationship between the input and the basis—retains most of the image’s semantic content. For instance, reconstructing an image using its original phase spectrum in combination with a random magnitude spectrum often results in an output that maintains the initial semantics. This observation indicates that altering neuron orientations is essential for semantic modifications to the generated images—a central objective in both subject-driven and controllable generation tasks. However, granting too much flexibility in adjusting neuron directions risks degrading the performance obtained during pretraining. To regulate this flexibility, we advocate for the preservation of pairwise neuron angles, positing that these angular relationships are critical for encoding neural network knowledge. Guided by this reasoning, we introduce a method to learn a layer-wise, shared orthogonal transformation for neurons within each layer so that hyperspherical energy remains constant.

    \item We introduce a new finetuning technique called Orthogonal Finetuning, designed to enhance the controllability of text-to-image diffusion models. To address stability concerns, we further devise a constrained version that restricts the angular shift from the original pretrained model.
    \item In contrast with existing finetuning strategies, OFT enables model adaptation while rigorously maintaining hyperspherical energy—an attribute we empirically observe to correlate with the semantic integrity retained from the pretrained generator.
    \item We deploy OFT in two domains: subject-driven generation and controllable generation. Extensive experiments reveal that our approach notably outperforms previous methods in generation fidelity, convergence rate, and robustness during finetuning. Furthermore, OFT is highly data-efficient, achieving strong performance using substantially fewer training images and epochs.
    \item For assessing controllability, we propose a novel control consistency metric. This metric works by inferring the control signal from the generated output and measuring its alignment with the original input control. The results further substantiate the superior controllability offered by OFT.
\end{itemize}

Drawing on empirical findings that hyperspherical similarity serves as an effective semantic encoder~\cite{liu2017deep,liu2018decoupled,chen2020angular}, our methodology employs hyperspherical energy~\cite{liu2018learning} to represent the inter-neuron relational structure within a network layer. This quantity, hyperspherical energy, is calculated as the aggregated hyperspherical similarities (\emph{e.g.}, cosine similarity) among all neuron pairs in a given layer, thereby indicating how evenly distributed neurons are on the surface of a unit hypersphere~\cite{liu2021learning}. We conjecture that an optimally finetuned model should exhibit a minimal shift in hyperspherical energy relative to the original pretrained model. While it is straightforward to regularize finetuning by penalizing any change in hyperspherical energy, this strategy alone does not reliably reduce the actual difference. Instead, we exploit a unique invariance property: the pairwise hyperspherical similarity remains constant under any shared orthogonal transformation of the neuron set. This motivates our introduction of Orthogonal Finetuning~(OFT), a strategy to adapt large text-to-image diffusion models to new downstream tasks while preserving the original hyperspherical energy. OFT operates by learning a single orthogonal transformation, shared across layers, that preserves all pairwise angles between neurons, effectively reorienting the canonical basis of the layer without distorting its internal geometry. Furthermore, recognizing that a smaller Euclidean distance between the finetuned and pretrained models helps maintain pretraining advantages, we develop a variant called Constrained Orthogonal Finetuning~(COFT). COFT restricts the adaptation process such that finetuned neurons are constrained within a hypersphere of constant radius centered at the pretrained configuration.

\textbf{Controllable generation}. Image-to-image translation represents a specialized case of controllable generation. Traditional approaches often utilize conditional generative adversarial networks~\cite{isola2017image,zhu2017toward,wang2018high,park2019semantic,choi2018stargan} for this purpose. In addition, diffusion models have been explored for image-to-image translation~\cite{saharia2022palette,voynov2022sketch,wang2022pretraining}. Recent advances such as ControlNet~\cite{zhang2023adding} introduce techniques to exert control over a pretrained diffusion model by finetuning it to respond to supplementary control inputs, resulting in remarkable controllable generation outcomes. Similarly, the T2I-Adapter~\cite{mou2023t2i} independently follows this strategy by finetuning pretrained diffusion models to enhance controllability in generated imagery. Consistent with the frameworks of \cite{zhang2023adding,mou2023t2i}, our method utilizes OFT for finetuning diffusion models, consistently achieving superior controllability using fewer training examples and requiring less finetuning of model parameters. Importantly, OFT does not impose additional computational burden during inference at test time.

Additionally, our approach is informed by orthogonal over-parameterized training~\cite{liu2021orthogonal}, which initiates classification neural networks from random initializations and orthogonally transforms them during training. The underlying principle is that random initializations yield expectedly low hyperspherical energy, and the objective of \cite{liu2021orthogonal} is to maintain this low energy throughout training—a property empirically linked to enhanced classification generalization~\cite{liu2018learning,lin2020regularizing}. Their findings demonstrate that orthogonal transformation possesses enough capacity for effective training of generalizable classifiers. In contrast, our primary concern is adapting text-to-image diffusion models to enhance controllability and improve generative outcomes on downstream tasks. We highlight two important distinctions between our method, OFT, and \cite{liu2021orthogonal}. Firstly, while the latter seeks to minimize hyperspherical energy, OFT is specifically designed to retain the hyperspherical energy of the original pretrained model, safeguarding the pretrained structure during finetuning—a crucial consideration since reducing hyperspherical energy during diffusion model adaptation may undermine semantic coherence. Secondly, OFT includes an explicit objective to reduce the deviation from the pretrained model, thereby producing a constrained variant, whereas \cite{liu2021orthogonal} does not employ such restrictions. We contend that the central challenge of finetuning lies in balancing adaptability and retention of pretrained knowledge, a balance that OFT is expressly constructed to achieve. Our main contributions are summarized as follows:

\textbf{Experimental details}. DreamBooth~\cite{ruiz2023dreambooth} and LoRA~\cite{hulora2022} serve as the baseline methods. All techniques utilize the loss function described in DreamBooth. For both DreamBooth and LoRA, we closely adhere to their respective methodologies and utilize their optimal hyperparameter configurations. Additional results and methodological specifics can be found in Appendix~\ref{app:settings},\ref{app:coft_oft},\ref{app:more_qualitative},\ref{app:failure}.

To provide further intuition, take a fully connected layer $\thickmuskip=2mu \medmuskip=2mu \bm{W}=\{\bm{w}_1,\cdots,\bm{w}_n\}\in\mathbb{R}^{d\times n}$, where $\bm{w}_i\in\mathbb{R}^{d}$ is the $i$th neuron and $\bm{W}^0$ denotes its pretrained weights. The layer's output vector $\thickmuskip=2mu \medmuskip=2mu \bm{z}\in\mathbb{R}^n$ is given by $\thickmuskip=2mu \medmuskip=2mu \bm{z}=\bm{W}^\top\bm{x}$ for an input $\thickmuskip=2mu \medmuskip=2mu \bm{x}\in\mathbb{R}^d$. The OFT approach seeks to minimize the difference in hyperspherical energy between the updated and original model:

\textbf{Connection to LoRA}. LoRA introduces a low-rank matrix to mitigate drastic changes in the pretrained weight matrix, thereby preserving the original model information. In comparison, OFT regulates the transformation of the pretrained weights via an orthogonal (full-rank) matrix, ensuring that the transformation does not compromise pretraining knowledge. The forward computation in OFT can be expressed as $\thickmuskip=2mu \medmuskip=2mu \bm{z}=(\bm{R}\bm{W}^0)^\top\bm{x}=(\bm{W}^0+(\bm{R}-\bm{I})\bm{W}^0)^\top\bm{x}$, with $\thickmuskip=2mu \medmuskip=2mu (\bm{R}-\bm{I})\bm{W}^0$ functioning analogously to LoRA’s low-rank adjustment. Given that $\bm{W}^0$ generally has full rank, OFT can also yield a low-rank adaptation if $\thickmuskip=2mu \medmuskip=2mu \bm{R}-\bm{I}$ possesses low rank. Both methods introduce a parameter controlling the number of trainable components: LoRA utilizes a rank parameter $r'$, while OFT makes use of a block-diagonal parameter $r$ for parameter reduction. Fundamentally, LoRA and OFT exemplify two different parameter-efficient mechanisms: LoRA leverages low-rank representations to minimize learnable parameters, whereas OFT capitalizes on sparsity via block-diagonal orthogonality.

\textbf{No extra inference cost}. In contrast to ControlNet, the OFT approach incurs zero inference-time overhead on the modified model. During inference, one can directly integrate the trained orthogonal matrix $\bm{R}$ with the initial weight matrix $\bm{W}^0$, resulting in an updated weight matrix $\thickmuskip=2mu \medmuskip=2mu \bm{W}=\bm{R}\bm{W}^0$. This procedure preserves the inference efficiency of the original pretrained model.

\textbf{Why OFT converges faster}. As demonstrated in Figure~\ref{fig:toy_ae}, the most impactful way to alter semantic representations is through neuron direction changes, precisely the aim of OFT. Furthermore, OFT's approach can be interpreted as updating neuron parameters constrained to a smooth hypersphere manifold, which offers more favorable optimization dynamics. Empirical support for this view is also found in \cite{liu2021orthogonal}.

\textbf{Text-to-image diffusion models}. Significant advancements~\cite{saharia2022photorealistic,ramesh2022hierarchical,rombach2022high,gu2022vector,nichol2022glide} in text-to-image synthesis are largely driven by breakthroughs in diffusion-based generative modeling~\cite{ho2020denoising,song2021score,song2021denoising,dhariwal2021diffusion} and integrated vision-language learning~\cite{su2020vlbert,lu2019vilbert,radford2021learning,wang2021simvlm,li2022blip,li2023blip,alayrac2022flamingo,schuhmann2022laion}. GLIDE~\cite{nichol2022glide} and Imagen~\cite{saharia2022photorealistic} operate in the pixel domain: GLIDE aligns its text encoder with a diffusion prior using text-image pairs, whereas Imagen leverages a frozen pretrained text encoder. Latent-space diffusion is explored by Stable Diffusion~\cite{rombach2022high} and DALL-E2~\cite{ramesh2022hierarchical}, where Stable Diffusion adopts VQ-GAN~\cite{esser2021taming} to build its latent visual vocabulary and DALL-E2 employs CLIP~\cite{radford2021learning} for constructing a unified image-text embedding. Apart from diffusion techniques, text-to-image synthesis has also utilized generative adversarial networks~\cite{reed2016generative,zhang2017stackgan,xu2018attngan,li2019controllable} and autoregressive models~\cite{ramesh2021zero,ding2021cogview,wu2022nuwa,yuscaling}. Notably, OFT provides a general-purpose finetuning framework applicable to any text-to-image diffusion model.

\textbf{Qualitative comparison}. We further provide a qualitative assessment of OFT and COFT alongside leading methods such as ControlNet, T2I-Adapter, and LoRA. As illustrated by the randomly generated samples in Figure~\ref{fig:control_final}, both OFT and COFT consistently produce images with high realism and fidelity, while maintaining precise controllability. For the S2I task, it is evident that LoRA fails to adhere to the input segmentation maps, whereas ControlNet, OFT, and COFT succeed in accurately steering the output images according to the provided controls. Notably, OFT and COFT surpass ControlNet by synthesizing more visually detailed and geometrically coherent images, despite using substantially fewer model parameters. In the C2I scenario, OFT and COFT can generate plausible images based on basic Canny edges, outperforming both T2I-Adapter and LoRA. Regarding the L2F task, our approach achieves the highest accuracy in pose control, even under difficult facial orientations. Across all three control scenarios, OFT and COFT deliver superior image quality compared to state-of-the-art baselines, underscoring the robustness of the OFT framework for controllable image synthesis. To enhance this qualitative evaluation, we supply additional visual comparisons for all control tasks in Appendix~\ref{app:control_exp}, and further demonstrate that OFT extends effectively to various other control scenarios (including dense pose-to-human body, sketch-to-image, and depth-to-image) in Appendix~\ref{app:more_control_task}.

We next examine parameter efficiency by comparing OFT with LoRA. In LoRA, for a low-rank value $r'$, the number of trainable parameters is given by $\thickmuskip=2mu \medmuskip=2mu r'(d+n)$. Assuming $r$ and $r'$ scale with the model dimension $d$, i.e., $\thickmuskip=2mu \medmuskip=2mu r=r'=\alpha d$ for some constant $0<\alpha\leq1$, the parameter complexity for LoRA is $\mathcal{O}(d^2+dn)$, while for OFT it becomes $\mathcal{O}(d)$. For illustrative purposes, consider a $128\times 128$ weight matrix; LoRA with $\thickmuskip=2mu \medmuskip=2mu r'=8$ introduces $2,048$ parameters, whereas OFT with $\thickmuskip=2mu \medmuskip=2mu r=8$ (and no block sharing) only introduces $960$ trainable parameters.

\begin{equation}\label{eq:oft_general}
    \footnotesize
    \bm{z}=\bm{W}^\top\bm{x}=(\bm{R}\cdot\bm{W}^0)^\top\bm{x},~~~\text{s.t.}~\bm{R}^\top\bm{R}=\bm{R}\bm{R}^\top=\bm{I}
\end{equation}

Traditional orthogonalization approaches such as the Gram-Schmidt process are differentiable but typically computationally prohibitive for large-scale applications~\cite{liu2021orthogonal}. To improve computational efficiency, we utilize the Cayley parameterization to generate the orthogonal matrix. This approach constructs $\bm{R}$ according to $\thickmuskip=2mu \medmuskip=2mu \bm{R}=(\bm{I}+\bm{Q})(I-\bm{Q})^{-1}$, where $\bm{Q}$ is skew-symmetric, satisfying $\thickmuskip=2mu \medmuskip=2mu \bm{Q}=-\bm{Q}^\top$. While this method greatly reduces computational costs, it imposes the minor restriction that only orthogonal matrices with determinant 1 (elements of the special orthogonal group) can be produced. Our experiments indicate this constraint does not negatively impact empirical results. Nevertheless, direct parameterization of $\bm{R}$ through the Cayley transform can still result in parameter inefficiency when $d$ is large. To alleviate this, we introduce a block-diagonal structure for $\bm{R}$, dividing it into $r$ blocks, as expressed below:

\section{Introduction}

\begin{itemize}[leftmargin=*,nosep]

We begin by articulating the rationale for employing orthogonal transformations in finetuning text-to-image diffusion models. This inquiry can be divided into two central aspects: (1) the motivation behind adjusting neuron angles (\emph{i.e.}, modifying direction), and (2) the justification for using orthogonal transformations specifically to effect angle finetuning.

At their core, both tasks focus on the challenge of how to finetune text-to-image diffusion models in a manner that maintains the strong generative capabilities achieved during pretraining. The primary criteria for an effective finetuning approach can be delineated as follows: (1) \emph{training efficiency}, which entails minimizing the number of parameters requiring updates and reducing the total number of training epochs, and (2) \emph{generalizability preservation}, which concerns retaining the model’s ability to generate outputs that are both high-quality and diverse. Typical finetuning methodologies fall into two broad categories: modifying existing neuron weights incrementally with a small learning rate (\emph{e.g.}, \cite{ruiz2023dreambooth}), or augmenting the model architecture with compact modules featuring re-parameterized neuron weights (\emph{e.g.}, \cite{hulora2022,zhang2023adding}). However, while these techniques offer simplicity, neither provides assurances that the model’s original generative performance will be preserved. Moreover, the field currently lacks rigorous theoretical guidelines for constructing effective finetuning strategies or for setting hyperparameters such as the number of training epochs. A central challenge stems from the absence of an explicit metric that evaluates how well pretrained generative abilities are retained. Prevailing finetuning practices operate under the assumption that minimizing the Euclidean distance between the parameters of the finetuned and pretrained models correlates with superior preservation of generative functionality. Consequently, finetuning is often conducted with exceptionally low learning rates. Yet, while this heuristic can be effective in certain scenarios, we contend that relying solely on Euclidean proximity is insufficient for assessing semantic retention. Thus, a more principled and structural approach for comparing finetuned and pretrained models could significantly enhance both the stability of finetuning and the fidelity of pretraining performance retention.

\textbf{Experimental protocols}. Our experimental setup is based on the pretrained text-to-image model Stable Diffusion v1.5~\cite{rombach2022high}. To maintain objectivity, we sample images randomly from the outputs of each method. Subject-driven generation experiments follow the approach outlined by DreamBooth~\cite{ruiz2023dreambooth}, while experiments on controllable generation follow the methodologies of ControlNet~\cite{zhang2023adding} and T2I-Adapter~\cite{mou2023t2i}. To ensure parity with LoRA, OFT or COFT is employed only at those layers where LoRA is applied. Comprehensive details and further results are documented in Appendix~\ref{app:settings}.

\section{Intriguing Insights and Discussions}

Recent advances in text-to-image diffusion models~\cite{saharia2022photorealistic,ramesh2022hierarchical,rombach2022high} have yielded strong results in producing high-quality images under textual supervision. Nevertheless, relying solely on text input often results in ambiguities and fails to afford the level of precise, fine-grained control required for certain applications. In this work, we focus on two specific tasks within the domain of text-to-image generation:

\subsection{Why Does Orthogonal Transformation Make Sense?}

\begin{equation}\label{eq:oft_principle}
\footnotesize
\min_{\bm{W}} \left\| \text{HE}(\bm{W})-\text{HE}(\bm{W}^0)\right\|~~\Leftrightarrow~~\min_{\bm{W}} \bigg{\|} \sum_{i\neq j} \|\hat{\bm{w}}_i-\hat{\bm{w}}_j\|^{-1}-\sum_{i\neq j} \|\hat{\bm{w}}^0_i-\hat{\bm{w}}^0_j\|^{-1}\bigg{\|}
\end{equation}

To further restrict the capacity of the standard OFT, one can confine the finetuned parameters to a neighborhood centered on the pretrained solution. The COFT variant applies the following transformation during the forward pass:

\textbf{Subject-driven generation}. To safeguard subject attributes,~\cite{avrahami2022blended,nichol2022glide} integrate user-provided masks as additional guidance. Inversion-based approaches~\cite{choi2021ilvr,dhariwal2021diffusion,ramesh2022hierarchical,gal2022image} allow modification of contextual features while preserving the core subject. Other methods, such as~\cite{hertz2022prompt}, enable both local and global editing without necessitating external masks, although these may not fully retain identity-specific features. Pivotal Tuning~\cite{roich2022pivotal} adapts the generator near an initially inverted code, with regularization to maintain identity, and~\cite{nitzan2022mystyle} learns a subject-specific face prior from a dataset of individual images. Instance-level variations are produced in~\cite{casanova2021instance}, but at the potential cost of losing fine subject details. DreamBooth~\cite{ruiz2023dreambooth} introduces a customized token and uses limited subject images to optimize the text-to-image diffusion model through reconstruction and class-specific prior preservation losses. Our method, OFT, operates within the DreamBooth framework but distinguishes itself by employing orthogonal transformations rather than standard low-rate finetuning.

Driven by the insight that angular relationships among neurons are essential for defining visual semantics, this work introduces a straightforward yet powerful finetuning technique—orthogonal finetuning—to enhance control over text-to-image diffusion models. The approach specifically addresses two key text-to-image scenarios: subject-driven generation and controllable generation. When evaluated against prior methods, OFT achieves superior controllability and greater finetuning stability, all while requiring fewer finetuning parameters. Notably, OFT imposes no additional computation during inference, making it a practical choice for deployment.

\textbf{OFT is adaptable across diverse architectures.} In theory, the OFT approach is applicable to any neural network architecture. In Transformer models, methods such as LoRA are commonly applied to the attention matrices~\cite{hulora2022}. For a fair comparison with LoRA in our evaluations, we limit OFT modifications to the attention layers. OFT's framework also lends itself well to convolutional networks; unlike LoRA, the block-diagonal configuration of $\bm{R}$ carries interpretable meaning in the convolutional context. Matching the number of $\bm{R}$'s blocks to the input channels results in each block transforming a specific channel, akin to a depth-wise convolution~\cite{chollet2017xception}. When all $\bm{R}$ blocks are identical and shared, OFT applies a universal orthogonal transform across all channels. Further investigations into OFT's application to convolutional finetuning are provided in Appendix~\ref{app:convolution}.

\textbf{Qualitative comparison}. To provide a clearer, more intuitive insight into the advantages offered by OFT, we display a selection of randomly chosen samples of subject-driven image generation in Figure~\ref{fig:dreambooth_final}. To ensure fairness in our comparison, every example is generated using the same finetuned model under each method, with no prompt-specific optimization carried out for any final output. For each approach, we utilize the model that records the highest CLIP validation metrics. Reviewing the samples in Figure~\ref{fig:dreambooth_final}, it is evident that both OFT and COFT exhibit strong semantic fidelity to the intended subject, whereas LoRA frequently fails to maintain subject identity (\emph{e.g.}, the subject identity is entirely lost by LoRA in the bowl example). Furthermore, OFT and COFT offer notably more precise alignment with the given text prompts, contrasting with DreamBooth, which—despite effectively retaining the subject—often produces images that do not align with the text (\emph{e.g.}, see the first row of the bowl sample). This qualitative analysis demonstrates that our OFT method provides superior controllability in addition to robust subject preservation. Additionally, OFT's performance remains stable irrespective of the number of training iterations, successfully scaling to more iterations—a property not shared by DreamBooth or LoRA. Further qualitative results are presented in Appendix~\ref{app:more_qualitative}. We also perform a human evaluation study described in Appendix~\ref{app:human}, which further confirms OFT’s performance benefits.

\section{Orthogonal Finetuning}

\textbf{Quantitative comparison}. To quantitatively assess model performance, following the protocol in \cite{ruiz2023dreambooth}, we measure subject fidelity (DINO~\cite{caron2021emerging}, CLIP-I~\cite{radford2021learning}), alignment to text prompts (CLIP-T~\cite{radford2021learning}), and sample diversity (LPIPS~\cite{zhang2018unreasonable}). Specifically, CLIP-I measures the mean pairwise cosine similarity between CLIP representations of generated and real images; DINO operates analogously, but employs ViT S/16 DINO features. The CLIP-T metric reflects the mean cosine similarity between text prompt embeddings and those of generated images. For diversity, we report the mean LPIPS cosine similarity among generated outputs conditioned on identical subjects and text prompts. As presented in Table~\ref{table:dreambooth_final}, both COFT and OFT substantially surpass DreamBooth and LoRA with respect to DINO and CLIP-I, and achieve competitive or superior outcomes on prompt fidelity and diversity. To mitigate the impact of stochasticity, each approach is evaluated by repeatedly finetuning the same pretrained model across 30 unique random seeds. Overall, our results confirm that the OFT framework not only offers enhanced convergence and stability, but also produces improved and consistently superior final results.

\section{Related Work}

\textbf{Model finetuning}. The strategy of finetuning large pretrained models to adapt to new tasks has become increasingly widespread~\cite{devlin2018bert,brown2020language,he2022masked}. Adaptation-based finetuning techniques (such as those presented in \cite{hulora2022,houlsby2019parameter,pfeiffer2020adapterfusion}) have been widely investigated within the realm of natural language processing. LoRA~\cite{hulora2022}, which assumes an additive low-rank decomposition for weight updates in finetuning, shares close relevance with OFT. In distinction, OFT applies a multiplicative, layer-shared orthogonal transformation to modify neuron weights, guaranteeing preservation of the pairwise angles between neurons within a given layer, thereby improving stability.

\subsection{Subject-driven Generation}

\subsection{Constrained Orthogonal Finetuning}

\subsection{Controllable Generation}

\begin{equation}\label{eq:coft_general}
    \footnotesize
    \bm{z}=\bm{W}^\top\bm{x}=(\bm{R}\cdot\bm{W}^0)^\top\bm{x},~~~\text{s.t.}~\bm{R}^\top\bm{R}=\bm{R}\bm{R}^\top=\bm{I},~\norm{\bm{R}-\bm{I}}\leq \epsilon
\end{equation}

\textbf{Finetuning stability and convergence}. We begin by assessing both the stability and convergence rate during finetuning across DreamBooth, LoRA, OFT, and COFT, as depicted in Figure~\ref{fig:teaser} and Figure~\ref{fig:db_convergence}. Our analysis indicates that COFT and OFT provide stable finetuning behavior for the diffusion model. After completing 400 iterations, both DreamBooth and OFT-based approaches attain effective control over the output, whereas LoRA does not maintain subject identity adequately. With 2000 iterations, DreamBooth exhibits mode collapse in its generated samples, and LoRA struggles to render the correct attribute (yellow shirt), mistakenly generating yellow fur instead. By contrast, OFT and COFT maintain robust and consistent controllability of the generated content. These findings substantiate that our OFT approach achieves rapid yet robust convergence in subject-driven generation tasks. Importantly, our results reveal that the methods' resilience to the choice of iteration count is significant, as it obviates the need to individually adjust the number of finetuning steps for various subjects. For OFT and COFT, a relatively large number of iterations can be fixed in advance without requiring meticulous adjustment. Furthermore, COFT, with a suitable $\epsilon$ parameter, renders both the learning rate and iteration number essentially tuneless.

    \item \textbf{Subject-driven generation}~\cite{ruiz2023dreambooth}: Here, the challenge is to produce new images of a particular subject in varied scenarios, using only a handful of images as references along with a descriptive prompt.
    \item \textbf{Controllable generation}~\cite{zhang2023adding,mou2023t2i}: This task involves utilizing additional control information (such as canny edges or segmentation maps) to guide the image synthesis process in conjunction with the provided text prompt.
\end{itemize}

\section{Concluding Remarks and Open Problems}

\textbf{Settings}. For the baseline methods, we compare with ControlNet~\cite{zhang2023adding}, T2I-Adapter~\cite{mou2023t2i}, and LoRA~\cite{hulora2022}. Our main experiments span three challenging controllable generation scenarios: transforming Canny edges into images~(C2I) on COCO~\cite{lin2014microsoft}, converting segmentation maps to images~(S2I) on ADE20K~\cite{zhou2017scene}, and mapping landmarks to facial images~(L2F) on CelebA-HQ~\cite{karras2018progressive,xia2021tedigan}. All evaluated approaches are used to finetune Stable Diffusion~(SD) v1.5 across the three datasets, training for 20 epochs each. Supplementary results can be found in Appendix~\ref{app:more_qualitative}, \ref{app:more_control_task}, and \ref{app:failure}.

We introduce a straightforward augmentation to OFT that allows each neuron to possess its own trainable scaling factor. This enhancement is motivated by the observation that scaling neuron activations does not alter the hyperspherical energy, due to the normalization of magnitude. Concretely, the new forward computation is given by $\thickmuskip=2mu \medmuskip=2mu\bm{z}=(\bm{R}\bm{W}^0\bm{D})^\top\bm{x}$\footnote{\textbf{Errata}: In the NeurIPS camera-ready manuscript, the re-scaled OFT forward computation was mistakenly reported as $\thickmuskip=2mu \medmuskip=2mu\bm{z}=(\bm{D}\bm{R}\bm{W}^0)^\top\bm{x}$. However, the initial implementation is accurate; therefore, the corresponding results in Appendix~\ref{app:rescaled} are correct.}, with $\thickmuskip=2mu \medmuskip=2mu \bm{D}=\text{diag}(s_1,\cdots,s_n)\in\mathbb{R}^{n\times n}$ being a diagonal matrix whose learnable entries $s_1,\cdots,s_n$ are all positive. Unlike the original OFT forward pass in Eq.~\eqref{eq:oft_general}, where only $\bm{R}$ was updated during training, both $\bm{D}$ and $\bm{R}$ are optimized here. This modification marginally increases the parameter count, yet substantially extends OFT's adaptability. For experimental clarity and to isolate the effect of orthogonal transformations, we adhere to standard OFT in the main text, though re-scaled OFT is typically superior (refer to Appendix~\ref{app:rescaled} for results).

Traditionally, finetuning relies on gradient descent with a reduced learning rate to modify a model’s parameters (or only certain layers). The choice of a small learning rate acts as a soft regularization, implicitly encouraging minimal changes from the initial pretrained weights. In effect, classical finetuning can be viewed as training the model while simultaneously keeping $\thickmuskip=2mu \medmuskip=2mu \| \bm{M} - \bm{M}^0\|$ small, where $\bm{M}$ and $\bm{M}^0$ represent the finetuned and pretrained weights, respectively. Although this form of implicit regularization is reasonable, it may lack sufficient restrictiveness for tuning very large models. To mitigate this, LoRA enforces an explicit low-rank constraint on the updates: $\thickmuskip=2mu \medmuskip=2mu \text{rank}(\bm{M} - \bm{M}^0)=r'$, where $r'$ is chosen to be small. In contrast, OFT imposes a restriction on the pairwise relationships among neurons, specifically $\thickmuskip=2mu \medmuskip=2mu \| \text{HE}(\bm{M})-\text{HE}(\bm{M}^0) \|=0$, where $\text{HE}(\cdot)$ indicates the hyperspherical energy for a weight matrix. 

Turning to the second question, the most straightforward strategy for adjusting neuron directions is to apply a uniform rotation or reflection across all neurons within a layer, inherently implementing an orthogonal transformation. Although other angular transformations—such as individually rotating neurons by varying angles—offer additional flexibility, we find that orthogonal transformations strike an optimal balance between adaptability and structural consistency. Additionally, as discussed in \cite{liu2021orthogonal}, orthogonal transformations are sufficiently expressive for learning robust neural representations. To corroborate this perspective, we conducted an experiment assessing the regularization effects imparted by imposing orthogonality. Utilizing the experimental configuration from ControlNet~\cite{zhang2023adding}, the results, summarized in Figure~\ref{fig:ortho}, show that our typical OFT method operates reliably and yields precise control by training epoch 20. In contrast, omitting the orthogonality constraint from OFT results in poor generation performance with no effective control. These findings affirm that the orthogonality constraint is critical to the success of OFT.

\textbf{Why not minimize hyperspherical energy}. Unlike \cite{liu2021orthogonal}, our method does not attempt to minimize hyperspherical energy. In the context of classification, a set of neurons with minimal redundancy is beneficial. Achieving minimum hyperspherical energy disperses neurons evenly over the hypersphere, an objective that is ill-suited for finetuning tasks as it risks undermining pretrained model information.

\subsection{Re-scaled Orthogonal Finetuning}

\textbf{Convergence}. We assess the convergence rates of ControlNet, T2I-Adapter, LoRA, and COFT on the L2F task through both quantitative and qualitative analysis. For the quantitative metric, we calculate the average $\ell_2$ distance between the ground truth face landmarks used for control and the model-generated face landmarks. Figure~\ref{fig:face_convergence} illustrates the evolution of the landmark error as a function of the number of finetuning epochs for each approach. It is clear from these results that COFT and OFT attain faster convergence than the other methods. Specifically, LoRA requires 20 epochs to reach the accuracy that our OFT model attains after just 8 epochs. Notably, both OFT and COFT maintain a comparable count of trainable parameters to LoRA, which is substantially fewer than ControlNet, yet they converge more rapidly than prior techniques. Furthermore, OFT's accelerated convergence is also supported by the outcomes depicted in Figure~\ref{fig:teaser}, where the rightmost sample highlights OFT's superior data efficiency over ControlNet and LoRA; OFT achieves satisfactory convergence using merely 5\% of the ADE20K dataset. For the qualitative assessment, our focus is on OFT, COFT, and ControlNet, given that ControlNet produces a landmark error similar to ours. Figure~\ref{fig:face_convergence_qual} reveals that OFT and COFT consistently exhibit stable convergence, with generated facial poses progressively aligning to the control landmarks. In comparison, ControlNet shows a delayed emergence of controllability, only appearing after the 8-th epoch—mirroring the abrupt convergence observed in \cite{zhang2023adding}. The steady and predictable convergence pattern of OFT makes it more practical for application, as its training process is smoother and more reliable, which in turn facilitates the optimization of its hyperparameters.